
\documentclass{article}
\usepackage[center]{titlesec}
\usepackage{graphicx} % Required for inserting images
\usepackage{tikz}
\usepackage{tikz-cd}
\usepackage{amsmath}
\usepackage{amssymb}
\usepackage{amsthm}
\usepackage{comment}
\usetikzlibrary{cd}
\usepackage{lmodern}
\usepackage{hyperref}
\usepackage{cleveref}
\hypersetup{
	colorlinks=true,   % sets links to be colored (false by default)
	linktoc=all,       % links both the section name and page number (options: 'section' or 'page')
	linkcolor=blue,    % internal links color (e.g., to sections, figures)
	citecolor=green,   % color for citation links
	urlcolor=cyan,     % color for URL links
}


% Add references.bib to your src
\usepackage[style=numeric]{biblatex} 

% Add bibliography
\addbibresource{references.bib} 

\numberwithin{equation}{section}

\theoremstyle{definition}
\newtheorem{thm}{Theorem}[section]
\newtheorem{lem}[thm]{Lemma}

\theoremstyle{definition}
\newtheorem{defn}[thm]{Definition}

\theoremstyle{remark}
\newtheorem{rem}[thm]{Remark}

\theoremstyle{remark}
\newtheorem{exa}[thm]{Example}

\titleformat{\section}[block]{\normalsize\filcenter}{\thesection}{1em}{\MakeUppercase}


\title{Notes on Potential Topics}
\author{Kate Birch}
\date{\today}

\begin{document}

\pagenumbering{gobble}
\maketitle

\tableofcontents

\newpage

\pagenumbering{arabic}

\section{Enriched Categories Quick Notes}

	\begin{defn}
		A \textit{monoidal category} is a sextuple $ \mathcal{V}= ( \mathbb{V}, \otimes, I, a, l , r)$ where $ \mathbb{V}$ is a category, $ \otimes : \mathbb{V} \times \mathbb{V} \longrightarrow \mathbb{V}$ a functor, $I$ an object of $ \mathbb{V}$ and natural isomorphisms $ a _{XYZ}: (X\otimes Y) \otimes Z \longrightarrow X \otimes (Y\otimes Z)$, $ l_X:I\otimes X \longrightarrow  X$, $r_X: X\otimes I \longrightarrow  X$ such that the following two diagrams commute: \\

		These diagrams will be added later, but they just express that $ a$ is associative and $ r$ and $ l$ act like identities.
	\end{defn}
	
	\begin{defn}
		A \textit{cartesian monoidal category} is a special case of a monoidal category in which $ \mathbb{V}$ is at category with finite products, and $ \otimes, I$ are respectively $ \times$ and the terminal object $1$, with $ a,l,r$ being the canonical isomorphisms.
	\end{defn}

	\begin{exa}
		Examples of cartesian monoidal categories are $Set, Cat, $, $Groupoid, Ord, Top$.	
		An example of a monoidal category is the category of endofunctors for a small category, with $ \otimes$ being composition.

	\end{exa}
	
	\begin{defn}
		In a monoidal category $ \mathcal{V}$, we have the \textit{representable functor} $V= \hom _{ \mathbb{V}}(I, X): \mathbb{V} \longrightarrow Set$.
	\end{defn}
		
	\begin{defn}
		A $ \mathcal{V}$\textit{-category} is a category $ \mathcal{A}$ in which,  $ \hom _{ \mathcal{A}}(A,B)$ is an object in $ \mathcal{V}$, als
	\end{defn}
	

\section{Biadjoint Triangles}

	\begin{defn} 
		For an adjunction $ (F,U,\eta,\epsilon):\mathbb{X}\longrightarrow \mathbb{A}$, the associated $ comonad$ on $ \mathbb{A}$ is the triple $ S = (FU, \epsilon, F \eta U)$ where $ \epsilon: FU \longrightarrow 2_{ \mathbb{A}}$ and $ F \eta U: FU \longrightarrow FU^2$ are natural transformations satisfying the axioms dual to those of a monad.
	\end{defn}
		
	\begin{thm}
		Any comonad $ S = (FU, \epsilon, F \eta U)$ on $ \mathbb{A}$ gives rise to an adjoint pair $ (F_S,U_S,\eta_S,\epsilon_S):\mathbb{A}\longrightarrow \mathbb{A}_S$, where $ \mathbb{A}_S$ is the category of $ S$-coalgebras, $U_S(A)= (S(A), F \eta_{U(A)}) $ and $ F_S(A, h) = A$.
	\end{thm}
	
	\begin{thm}
		For an adjunction $ (F,U,\eta,\epsilon):\mathbb{X}\longrightarrow \mathbb{A}$ and associated comonad $ S = (FU, \epsilon, F\eta U)$ on $ \mathbb{A}$, there is a unique comparison functor $ K_S: \mathbb{X} \longrightarrow \mathbb{A}_S$ with $K_S U = U_S$ and $F_SK_S = F$. $ K$ is defined by $ X \longmapsto (F(X), F(\eta_X))$. 
	\end{thm}
	
	\begin{defn}
		%For an adjunction $ (F,U,\eta,\epsilon):\mathbb{X}\longrightarrow \mathbb{A}$, associated comonad  $ S = (FU, \epsilon, F\eta U)$ on $ \mathbb{A}$ and unique comparison functor $ K_S: \mathbb{X} \longrightarrow \mathbb{A}_S$, the functor $ F$ is \textit{precomonadic} if $K_:S$ is full and faithful.
		As in the situation above, the functor $ F$ is \textit{precomonadic} if $ K_S$ is full and faithful.
	\end{defn}

	\begin{defn}
		For a cocomplete, complete and symmetric monoidal closed category $ V$. Let $ (F \dashv U, \eta, \epsilon): \mathbb{X} \longrightarrow \mathbb{A} $ be a $ V$-adjunction.

		The associated $V$-natural isomorphism is written
		\[
		 \chi : \hom _{ \mathbb{A}}(F-,-) \simeq \hom _{ \mathbb{X}}(-,U-)
		\] 

		Where $ \chi_{(X,Z)}= \hom_{\mathbb{X}}(\eta_X, UZ) \circ U_{FX,Z}$.
	\end{defn}

		Note this is because $ U_{FX,Z}: \hom_{ \mathbb{A}}(FX,Z) \longrightarrow \hom_{ \mathbb{X}}(UFX, UZ) $ and $ \hom _{ \mathbb{X}}(\eta_X, UZ): hom_{ \mathbb{X}}(UFX, UZ) \longrightarrow \hom _{ \mathbb{X}}(X, UZ)$.

	\begin{thm}\label{thm:enrichedadjoint}
		\textit{Enriched Adjoint Triangle Theorem.} Let $ (F \dashv U, \eta, \epsilon)$ and $ (G \dashv V, \rho, \mu)$ be $ V$-adjunctions such that 		
			% https://q.uiver.app/#q=WzAsMyxbMiwwLCJcXG1hdGhiYntZfSJdLFsxLDEsIlxcbWF0aGJie0F9Il0sWzAsMCwiXFxtYXRoYmJ7WH0iXSxbMCwxLCJHIl0sWzIsMSwiRiIsMl0sWzIsMCwiSiJdXQ==
		\[\begin{tikzcd}
			{\mathbb{X}} && {\mathbb{Y}} \\
			& {\mathbb{A}}
			\arrow["J", from=1-1, to=1-3]
			\arrow["F"', from=1-1, to=2-2]
			\arrow["G", from=1-3, to=2-2]
		\end{tikzcd}\]	
		is a commutative triangle of $V$-functors. Then assume for each pair of objects $ X \in \mathbb{X}$ and $ Y \in \mathbb{Y}$, the induced diagram 
				% https://q.uiver.app/#q=WzAsMyxbMCwwLCJcXGhvbV97XFxtYXRoYmJ7WX19KEooWCksWSkiXSxbMSwwLCJcXGhvbV97XFxtYXRoYmJ7QX19KEYoWCksRyhZKSkiXSxbMywwLCJcXGhvbV97XFxtYXRoYmJ7QX19KEYoQSksIEdWRyhZKSkiXSxbMCwxLCJHX3tKKFgpLFl9Il0sWzEsMiwiR197SihYKSxWRyhZKX1cXGNpcmMgXFxjaGlfeyhKKFgpLCBHKFkpKX0iLDAseyJvZmZzZXQiOi0xfV0sWzEsMiwiXFxob21fe1xcbWF0aGJie0F9fShGKFgpLEcoXFxldGFfWSkpIiwyLHsib2Zmc2V0IjoxfV1d
		\[\begin{tikzcd}
			{\hom_{\mathbb{Y}}(J(X),Y)} & {\hom_{\mathbb{A}}(F(X),G(Y))} && {\hom_{\mathbb{A}}(F(X), GVG(Y))}
			\arrow["{G_{J(X),Y}}", from=1-1, to=1-2]
			\arrow["{G_{J(X),VG(Y)}\circ \chi_{(J(X), G(Y))}}", shift left, from=1-2, to=1-4]
			\arrow["{\hom_{\mathbb{A}}(F(X),G(\rho_Y))}"', shift right, from=1-2, to=1-4]
		\end{tikzcd}\]
		is an equaliser in $ V$. The $ V$-functor $J$ has a right $V$-adjoint $W$ if and only if, for each object $ Y$ in $ \mathbb{Y}$, the $V$-equaliser of 
			% https://q.uiver.app/#q=WzAsMixbMCwwLCJVRyhZKSJdLFszLDAsIlVHVkcoWSkiXSxbMCwxLCJVRyhWKFxcbXVfRyhZKSlcXGV0YV97SlVHKFkpfSlcXHJob197VUcoWSl9IiwwLHsib2Zmc2V0IjotMX1dLFswLDEsIlVHKFxcZXRhX1kpIiwyLHsib2Zmc2V0IjoxfV1d
		\[\begin{tikzcd}
			{UG(Y)} &&& {UGVG(Y)}
			\arrow["{UG(V(\epsilon_G(Y))\rho_{JUG(Y)})\eta_{UG(Y)}}", shift left, from=1-1, to=1-4]
			\arrow["{UG(\rho_Y)}"', shift right, from=1-1, to=1-4]
		\end{tikzcd}\]	

		exists in the $V$-category $ \mathbb{X}$. In this case, this equaliser gives the value of $ W(Y)$.

	\end{thm} \\
				
		\textit{Beck's Theorem} also gives us a corollary.

	\begin{thm}
				Let $ (F \dashv U, \eta, \epsilon)$ and $ (G \dashv V, \rho, \mu)$ be $ V$-adjunctions such that 		
				% https://q.uiver.app/#q=WzAsMyxbMiwwLCJcXG1hdGhiYntZfSJdLFsxLDEsIlxcbWF0aGJie0F9Il0sWzAsMCwiXFxtYXRoYmJ7WH0iXSxbMCwxLCJHIl0sWzIsMSwiRiIsMl0sWzIsMCwiSiJdXQ==
			\[\begin{tikzcd}
				{\mathbb{X}} && {\mathbb{Y}} \\
				& {\mathbb{A}}
				\arrow["J", from=1-1, to=1-3]
				\arrow["F"', from=1-1, to=2-2]
				\arrow["G", from=1-3, to=2-2]
			\end{tikzcd}\]	
		is a commutative triangle of $V$-functors and $ G$ is $V$-precomonadic. The $ V$-functor $ J$ has a right $V$-adjoint $G$ if and only if, for each object $Y$ in $ \mathbb{Y}$, the $V$-equaliser of
			\[\begin{tikzcd}
				{UG(Y)} &&& {UGVG(Y)}
				\arrow["{UG(V(\epsilon_G(Y))\rho_{JUG(Y)})\eta_{UG(Y)}}", shift left, from=1-1, to=1-4]
				\arrow["{UG(\rho_Y)}"', shift right, from=1-1, to=1-4]
			\end{tikzcd}\]	
		exists in the $V$-category of $ \mathbb{X}$.
		In this case, these equalisers give the value of the right adjoint $ W$.

	\end{thm}
		
		In Chapters 2-4, various definitions in enriched category theory are given: pseudofunctors, pseudonatural transformations, the Yoneda embedding for 2-functors, biadjoint functors, biadjunctions, local equivalences, modifications, the bicategorical Yoneda Lemma, and the weighted bilimit.
		Descent objects and descent diagrams and effective descent are then defined.
		Then \Cref{thm:enrichedadjoint} is given a more general formulation in terms of biadjunctions and descent objects (instead of equalisers).\\

		In Chapter 5, more definitions are given: pseudocomonads, pseudocoalgebras, and pseudoprecomonadicity. 
		Then \Cref{thm:enrichedadjoint} is given in terms of pseudoprecomonads.\\

		In Chapter 6, there is a discussion about the role of the unit and counit of the biadjunction in \Cref{thm:enrichedadjoint}.\\

		In Chapter 7, a characterisation of pseudocomonadicity is given in terms of creation of descent diagrams.\\

		In Chapter 8, there is a 2-(co)monadic approach to coherence, which is given by studying the inclusion of the 2-category of strict (co)algebras into the 2-category of pseudo(co)algebras of a given 2-(co)monad.\\

		In Chapter 9, a result of the \textit{Adjoint Triangle Theorem} is given. The result deals with the lifting of adjunctions to adjunctions between the Eilenberg-Moore categories. With this, a natural extension of the notion of Kan extensions is given for the tricategory $2-CAT$.
		Then the pointwise pseudo-Kan extension is discussed as well as conditions for its existence (amongst other results).


\section{On Lifting of Biadjoints and Lax Algebras}		


	\begin{defn}
		A \textit{2-category} is a category enriched over $Cat$ such that its collection of objects is a discrete category of $ CAT$. 
		This means, concretely, that in a 2-category, there are objects and morphisms (referred to respectively as 0- and 1-cells), but also morphisms between morphisms (called 2-cells), with appropriate vertical and horizontal composition.
		Given morphisms $ f,g:A \longrightarrow B$, we write $ \alpha: f \implies g$ to denote a 2-cell from $f $ to $g$.

	\end{defn}

	Category with topology on it or something
	Category that generalises shapes and geometry
	 
	\begin{defn}
		A \textit{weighted limit} over a functor $ F:K \longrightarrow \mathbb{C}$ (where $ K$ a small category) with the weight $W: K \longrightarrow \mathcal{V}$ for a $ \mathcal{V}$-enriched category $ \mathbb{C}$ is the object $ \lim ^W F$ in $ \mathbb{C}$ such that for all $C$ in $ \mathbb{C}$:
		\[
			\text{hom} _{ \mathbb{C}}(C, \lim ^W F) \simeq \text{hom} _{ \mathcal{V}^K}(W(-), \text{hom} _{ \mathbb{C}}(C, F(-)))
		\] 
	\end{defn}
	
	\begin{defn}
		Let $R: \mathbb{U} \longrightarrow  \mathbb{B}$and $ E: \mathbb{B} \longrightarrow \mathbb{U}$ be pseudofunctors. $ E$ is \textit{left biadjoint} to $ R$ if there exist
			\begin{itemize}
				\item pseudonatural transformations $ \rho: \text{Id} _{ \mathbb{B}} \longrightarrow RE $ and $ \epsilon: ER \longrightarrow \text{Id} _{ \mathbb{U}}$ \\
				\item invertible modifications $v: \text{id} _{E} \implies ( \epsilon E) (E \rho) $ and $ w: (R \epsilon)( \rho R) \implies \text{id} _{ \mathbb{ U}}$
			\end{itemize}
		satisfying coherence axioms.
	\end{defn}
	
	\begin{defn}
		Let $ R$ and $ E$ be as above, except now they are also 2-functors. Let $ (E \dashv R, \rho, \epsilon, v, w) $ be the biadjunction as described above. 
		Then $ E$ is \textit{left 2-adjoint} to $ R$ if $ v,w$ are identities and $ \rho, \epsilon$ are 2-natural transformations. 
		In this case, we write $ (E \dashv R, \rho, \epsilon)$ to be a 2-adjunction.
	\end{defn} 
	
	The paper introduces the definition of a computad, which can essentially be thought of as a graph that generates a category. 
	A small computad can be thought in the following way: as a graph (like "objects"). 
	Then we take the free category over this graph and take the underlying graph of that (so the same objects, but now we also have morphisms). This comes with a domain and codomain map into the computad.
	Lastly, a computad also has a way to combine morphisms. \\

	There is a generalisation of the biadjoint triangle theorem in terms of the weighted bicolimit, which turns out to define the right biadjoint we are interested in.\\

	In Section 3, lax descent objects are discussed. 
	I honestly found the definition very high level and difficult to appreciate. It is defined in terms of computads and weighted bilimits.
	I should read through Janelidze's Galois Theory notes.\\


	Section 4 describes pseudomonads, which are similarly defined as biadjunctions.
	The 2-category of \textit{lax $ \mathcal{T}$-algebras} is almost like normal $ \mathcal{T}$-algebras, (where now $ \mathcal{T}$ is a pseudomonad), but with 2-cells as well.
	There is also a sub-2-category of lax $ \mathcal{T}$-algebras, called pseudoalgebras, which I think is just the same thing, but without the 2-cells in the definition.
	Then the adjoint from this category is discussed.
	There is some more stuff about this.\\

	Lastly, we talk about conditions under which the inclusion of pseudo lax $ \mathcal{T}$-algebras into lax $ \mathcal{T} $-algebras is discussed (it turns out we need codescent objects in $Ps- \mathcal{T}-Alg$).

	These more summarised parts of the paper are more to do with the lifting of biadjoints.
	

		
		
	


\end{document}
